\chapter{Conclusion}

Enhancing board games using computer vision on mobile devices requires developers to fine-tune each application for the specific game. However, the foundations of these applications are mostly similar, each requiring a camera setup, detecting objects, loading assets, and storing user data. The goal of this thesis was to create a library containing these functionalities and providing a simple way for developers and hobbyists to create these applications. All of this was achieved while maintaining a modular approach and not forcing developers into a closed system that cannot be further fine-tuned.

The developed Kotlin library is built around Android's CameraX and OpenCV for image manipulation. Furthermore, splitting CVLiBG into two modules allows it to be relatively lightweight and to include Open Neural Network Exchange only as an option, should the developers choose to use it. The overall setup needed is relatively trivial as it mostly requires a few UI elements and building the camera controller alongside detectors.

During the development of the demonstration application and its dataset, several limitations of existing open-source annotation tools were encountered. These included slow workflows, instability, and the need to switch between multiple tools. This motivated the development of Dataset Creator as an all-in-one solution for annotating data, generating synthetic samples, and training object detection models inside a single application.

The developed application for the Bang! card game served as a proof of concept, demonstrating that CVLiBG can simplify the integration of object detection into Android applications. The library abstracts most camera and detection related tasks and provides developers with a clear interface for collecting and reacting to detection results. Furthermore, CVLiBG's services proved useful for repeatable tasks, such as loading image assets, storing game data, and saving user preferences.

Future extensions of the library could include other popular inference runtimes such as TensorFlow Lite, which would give developers more options for optimization on different hardware configurations. For the dataset creator, future extensions could include the generation of synthetic data, utilizing new generative methods or different models and training options.

Although CVLiBG and Dataset Creator have their limitations and are not the best in their class, they provide a complete workflow for Android applications utilizing object detection. Together, they cover the process from dataset preparation and model training to deployment inside an Android application and allow developers to focus on the user interface and interactions with the results of detected objects.

